\documentclass{situationist-preprint}

\title{The Contour Theory of Musical Expression}
\msid{AES-07-a}
\msfamily{AES-07}
\msdiscipline{Aesthetics}
\mstype{Principal manuscript}
\mskeywords{musical expression; contour theory; resemblance; arousal theory; persona theory; expressiveness; emotion}

\begin{document}
\maketitle

\begin{abstract}
Music is commonly described as sad, joyful or agitated, but a sequence of sounds has no feelings, and the sadness of a piece need not be felt by its composer, its performer or its listeners. We argue that music is nonetheless genuinely expressive, and that it expresses emotions by resembling the characteristic expressive behaviour of people who have them: their gait, posture, gesture and the contour of their voices. We state three conditions on an account of musical expressiveness: that expressiveness does not require anyone to feel the emotion expressed, that it is heard in the music, and that music's expressive range is selective, rich for emotions with a distinctive behavioural profile and poor for emotions individuated by their thoughts. We show that the contour theory meets all three, that the arousal and expression theories fail the first, and that the persona theory, though compatible with the data, depends on the resemblance the contour theory identifies. Empirical work on the cues shared by vocal expression and music performance supports the account.
\end{abstract}

\section{The problem}

We say that the slow movement of Beethoven's \emph{Eroica} symphony is grief-stricken and that the finale of his Seventh is exuberant, and listeners agree in such descriptions to a considerable degree. Yet emotions belong to sentient beings, and music is a structure of sounds. \citet{Hanslick1854} drew the conclusion that music does not express definite feelings at all: its content is tonally moving forms, and what it can represent of feeling is only its dynamic aspect, its speed and slowness, its swelling and dying away. The question is whether music is expressive, and if it is, how.

The answers in the literature fall into four groups. \emph{Expression theories} hold that music is expressive of an emotion when the composer, in making it, expressed an emotion he felt \citep{Tolstoy1898} or worked out and clarified one \citep{Collingwood1938}. \emph{Arousal theories} hold that music is expressive of an emotion when it disposes listeners to feel it, or a response related to it \citep{Matravers1998}. \emph{Persona theories} hold that to hear music as expressive is to hear it as the expression of the emotions of an imagined agent \citep{Levinson1996,Robinson2005}. \emph{Resemblance theories} hold that music is expressive because it resembles, in its dynamic character, something emotional. \citet{Langer1942} took the resemblance to be to the inner forms of feeling. \citet{Kivy1980} and \citet{SDavies1994} took it to be to the outward expression of emotion: Kivy compared the expressiveness of music to that of the face of a Saint Bernard dog, which looks sad whatever the dog feels, and Davies to that of a weeping willow, whose drooping shape presents the appearance of sadness. Kivy called his view the contour theory. We defend it.

\section{Framework}

\begin{definition}[Expressive behaviour]\label{d:behav}
The \emph{characteristic expressive behaviour} of an emotion $E$ is the pattern of bodily and vocal behaviour by which people typically show $E$: its characteristic gait, posture, gesture, facial configuration, and vocal properties such as pitch level, pitch range, loudness, rate and articulation.
\end{definition}

\begin{definition}[Contour resemblance]\label{d:res}
A passage of music $M$ \emph{contour-resembles} the expressive behaviour of $E$ if $M$ shares with that behaviour a perceptually salient profile of dynamic features: tempo with rate of movement, melodic contour and register with pitch contour and level of the voice, loudness with vocal intensity, articulation and rhythm with the smoothness or abruptness of movement.
\end{definition}

\begin{definition}[Contour theory]\label{d:ct}
The \emph{contour theory} holds that $M$ is expressive of $E$ if and only if $M$ contour-resembles the characteristic expressive behaviour of $E$, and this resemblance is heard by suitably experienced listeners as the appearance of $E$.
\end{definition}

The definition makes the resemblance one that is heard, not merely one that obtains. Sad music is slow, low in pitch, narrow in range, quiet and legato, and sad people move slowly, speak quietly and at a low pitch, and droop; joyful music is fast, high, wide-ranging and lively, and so are joyful people. The listener hears the music's movement as having the look, or rather the sound, of the behaviour.

Three conditions constrain any account.

\begin{principle}[No expresser]\label{p:noexp}
Music can be expressive of $E$ although no one, neither composer nor performer nor listener, feels $E$ in connection with it.
\end{principle}

\begin{principle}[Audibility]\label{p:aud}
The expressiveness of music is heard in the music: it is grasped quickly, by listeners without technical knowledge, and with substantial agreement among those familiar with the musical idiom.
\end{principle}

\begin{principle}[Selectivity]\label{p:sel}
Music without words or programme is reliably expressive of emotions with a distinctive behavioural profile, such as sadness, joy, tranquillity and agitation, and not of emotions individuated chiefly by their thoughts, such as jealousy, shame, hope or envy.
\end{principle}

Principle~\ref{p:noexp} is supported by familiar facts. Composers write sad music on commission and in good spirits; performers play it while thinking of their fingering; listeners enjoy it without becoming sad. Principle~\ref{p:aud} is supported by ordinary experience and by experimental findings of agreement on the basic emotions expressed in music \citep{JuslinLaukka2003}. Principle~\ref{p:sel} is Hanslick's observation, accepted by \citet{Kivy1980}: hope and jealousy differ in what they are about, not in how the hopeful and the jealous walk or speak, and music cannot by itself convey what an emotion is about.

\section{The argument}

\begin{lemma}\label{l:ct}
The contour theory satisfies Principles~\ref{p:noexp}--\ref{p:sel}.
\end{lemma}

\begin{proof}
\emph{No expresser.} Contour resemblance is a relation between the music's audible features and a type of behaviour. It holds whatever anyone feels, as a willow's resemblance to a drooping figure holds whatever the willow, or its observer, feels.

\emph{Audibility.} The relevant features, tempo, register, contour, loudness and articulation, are perceptible without technical knowledge, and the resemblance of their profile to the corresponding profile of expressive behaviour is grasped by the same capacities by which we recognize emotion in voices and movements, which are fast and widely shared.

\emph{Selectivity.} Music can resemble the expressive behaviour of an emotion only if the emotion has a characteristic expressive behaviour. Emotions with a distinctive behavioural profile have one; emotions individuated by their thoughts share their behaviour with other emotions, so music that resembles it resembles equally the behaviour of all of them, and is expressive at most of the more generic state that the behaviour expresses. The theory therefore predicts the selectivity recorded in Principle~\ref{p:sel}.
\end{proof}

\begin{lemma}\label{l:rivals}
The expression theory and the arousal theory violate Principle~\ref{p:noexp}.
\end{lemma}

\begin{proof}
The expression theory makes a passage expressive of $E$ only if its composer felt, or worked out, $E$ in composing it; a passage composed without that feeling is not expressive of $E$ on this theory, contrary to Principle~\ref{p:noexp}. The arousal theory makes a passage expressive of $E$ only if it disposes listeners to feel $E$ or a related response; a passage heard as sad by listeners whom it does not sadden, and who are not disposed to be saddened by it, is not expressive of sadness on this theory, contrary to Principle~\ref{p:noexp}. Sophisticated arousal theories weaken the response to a feeling related to $E$ \citep{Matravers1998}; but they must then explain why the listener's response is related to sadness, and the most natural explanation is that the listener hears the music as sad, which presupposes the expressiveness the theory sets out to explain.
\end{proof}

\begin{lemma}\label{l:persona}
The persona theory presupposes contour resemblance.
\end{lemma}

\begin{proof}
On the persona theory, a passage is expressive of $E$ if it is readily heard as the expression of $E$ by an imagined agent \citep{Levinson1996}. The theory is compatible with Principles~\ref{p:noexp}--\ref{p:sel}, since the imagined agent need not be anyone actual. But it requires an explanation of why a given passage is readily heard as the expression of sadness rather than of joy. The natural answer, and the one persona theorists generally give, appeals to the music's resemblance to human expressive gesture and utterance. That answer is the contour theory's account of expressiveness, and the persona is then an imaginative elaboration of a resemblance already heard.
\end{proof}

\begin{theorem}\label{t:main}
Music is expressive, and a passage of music is expressive of an emotion in virtue of its contour resemblance to the characteristic expressive behaviour of that emotion.
\end{theorem}

\begin{proof}
Music satisfies the descriptions of Principles~\ref{p:noexp}--\ref{p:sel}: it is heard, with agreement, as sad, joyful or agitated. By Lemma~\ref{l:ct}, the contour theory explains these facts. By Lemma~\ref{l:rivals}, the expression and arousal theories cannot; by Lemma~\ref{l:persona}, the persona theory explains them only by invoking contour resemblance. Langer's resemblance to inner forms of feeling is ruled out by Principle~\ref{p:sel}: hope and jealousy have inner dynamics as distinctive as those of sadness, so that if music resembled inner feeling, it should express them as readily, which it does not. Contour resemblance is therefore the basis of musical expressiveness, and since music has it, music is expressive.
\end{proof}

The empirical evidence fits Theorem~\ref{t:main}. In a review of over a hundred studies, \citet{JuslinLaukka2003} found that the acoustic cues by which emotions are communicated in vocal expression and in music performance largely coincide: sadness is conveyed in both by slow rate, low intensity, low pitch and little variability, and happiness by fast rate, high intensity, high pitch and more variability. The authors interpret the finding as evidence that musical expression draws on the code of vocal expression, which is the mechanism the contour theory postulates.

\section{Objections and replies}

\begin{objection}[Cheap resemblance]\label{o:cheap}
Everything resembles everything in some respect \citep{Goodman1970}. If expressiveness were resemblance, every passage would express every emotion.
\end{objection}

\begin{reply}
The contour theory does not appeal to resemblance in some respect or other. Definition~\ref{d:res} fixes the respects, which are the dynamic features by which we recognize emotion in people, and Definition~\ref{d:ct} requires the resemblance to be heard. In those respects a slow, quiet, descending passage resembles sad behaviour and does not resemble joyful behaviour. The constraint Goodman's point imposes, that resemblance be specified in respects, is one the theory meets.
\end{reply}

\begin{objection}[Willows]\label{o:willow}
A weeping willow resembles a drooping figure, but it does not express sadness. Resemblance therefore yields at most the appearance of an emotion, not its expression.
\end{objection}

\begin{reply}
The contour theory is a theory of expressiveness, the property of presenting the appearance of an emotion, and on that understanding the willow is expressive of sadness, as \citet{SDavies1994} holds. Music differs from the willow in being made and performed so as to have this appearance, and in being attended to for it, so that its expressiveness belongs to the content of the work and is a proper object of critical attention. Nothing more is needed for music to be expressive in the sense in which we call it sad. Whether an emotion is also expressed in the sense in which a person expresses one is a question about composers, not about music.
\end{reply}

\begin{objection}[Conventions]\label{o:conv}
Much of music's expressiveness is conventional. The association of the minor mode with sadness in Western tonal music resembles no behaviour.
\end{objection}

\begin{reply}
The theory claims that contour resemblance is the basis of musical expressiveness, not that every expressive feature is a resemblance. \citet{Kivy1980} himself allowed conventional associations, and argued that some of them may have their origin in resemblance. Conventional features supplement and modulate an expressiveness already grounded in contour, and it is the contour features that carry expression across idioms, as the agreement found for tempo, loudness and register suggests \citep{JuslinLaukka2003}.
\end{reply}

\begin{objection}[Felt responses]\label{o:felt}
Listeners do feel something in response to sad music, and often feel sad. A theory that makes no reference to listeners' feelings leaves out what makes musical expression matter.
\end{objection}

\begin{reply}
The contour theory is an account of what makes music expressive, not of what listeners feel. It is compatible with listeners' feeling sad, and helps explain why they do: \citet{SDavies1994} argues that responses of this kind are mirroring responses to perceived expressiveness, as one may be saddened by being among sad-looking people. The response depends on the expressiveness and cannot therefore constitute it, which is the point of Lemma~\ref{l:rivals}.
\end{reply}

\section{What the result does not establish}

The theory explains the expressiveness of music without words or programme. It does not show that music cannot be expressive of thought-dependent emotions in combination with a text, a title or a dramatic context, which supply what the emotion is about; it predicts that such expressiveness is due to the combination. It does not settle the value of musical expressiveness, which need not be the chief value of music, and it is compatible with Hanslick's view that the beauty of music lies in its forms. Finally, it does not decide whether the appearance of emotion heard in music is perceived directly or through imaginative projection; the companion manuscript AES-07-b examines the notion of resemblance at issue.

\section{Conclusion}

Music is expressive. A passage of music is sad, joyful or agitated in virtue of resembling, in tempo, contour, register, loudness and articulation, the way sad, joyful or agitated people move and speak. This explains why musical expressiveness requires no one to feel the emotion, why it is heard in the music, and why music expresses the emotions it does and not others. Rival theories either fail these facts or rely on the resemblance the contour theory identifies.

\begin{thebibliography}{99}\small
\bibitem[Collingwood(1938)]{Collingwood1938} Collingwood, R.~G. (1938). \emph{The Principles of Art}. Clarendon Press.
\bibitem[Davies(1994)]{SDavies1994} Davies, S. (1994). \emph{Musical Meaning and Expression}. Cornell University Press.
\bibitem[Goodman(1970)]{Goodman1970} Goodman, N. (1970). Seven strictures on similarity. In L. Foster and J.~W. Swanson (Eds.), \emph{Experience and Theory}. University of Massachusetts Press.
\bibitem[Hanslick(1854)]{Hanslick1854} Hanslick, E. (1854). \emph{On the Musically Beautiful}. Trans. G. Payzant. Hackett, 1986.
\bibitem[Juslin and Laukka(2003)]{JuslinLaukka2003} Juslin, P.~N. and Laukka, P. (2003). Communication of emotions in vocal expression and music performance: Different channels, same code? \emph{Psychological Bulletin}, 129(5), 770--814.
\bibitem[Kivy(1980)]{Kivy1980} Kivy, P. (1980). \emph{The Corded Shell: Reflections on Musical Expression}. Princeton University Press.
\bibitem[Langer(1942)]{Langer1942} Langer, S.~K. (1942). \emph{Philosophy in a New Key}. Harvard University Press.
\bibitem[Levinson(1996)]{Levinson1996} Levinson, J. (1996). Musical expressiveness. In \emph{The Pleasures of Aesthetics: Philosophical Essays}. Cornell University Press.
\bibitem[Matravers(1998)]{Matravers1998} Matravers, D. (1998). \emph{Art and Emotion}. Clarendon Press.
\bibitem[Robinson(2005)]{Robinson2005} Robinson, J. (2005). \emph{Deeper than Reason: Emotion and Its Role in Literature, Music, and Art}. Clarendon Press.
\bibitem[Tolstoy(1898)]{Tolstoy1898} Tolstoy, L. (1898). \emph{What Is Art?} Trans. A. Maude.
\end{thebibliography}

\end{document}
